
\documentclass[12pt,letterpaper]{article}

\usepackage[utf8]{inputenc}
\usepackage[T1]{fontenc}
\usepackage[spanish]{babel}
\usepackage{newtxtext,newtxmath}
\usepackage{geometry}
\usepackage{setspace}
\usepackage{enumitem}
\usepackage{hyperref}
\usepackage{microtype}

% -------------------------------------------------
% CONFIGURACIÓN
% -------------------------------------------------
\geometry{
    letterpaper,
    top=2.5cm,
    bottom=2.5cm,
    left=2.5cm,
    right=2.5cm
}

\setstretch{1}
\setlength{\parindent}{0pt}
\setlength{\parskip}{6pt}

% -------------------------------------------------
% DOCUMENTO
% -------------------------------------------------
\begin{document}

% -------------------------------------------------
% TÍTULO Y DATOS
% -------------------------------------------------

\begin{center}
    {\large\textbf{El arte de obtener información útil}}

    \vspace{0.3cm}
    G. D. Monzón Flores\\
    7690-20-20745 Universidad Mariano Gálvez\\[0.15cm]

    Seminario de Tecnología de Información\\
    gmonzonf1@miumg.edu.gt
\end{center}

% -------------------------------------------------
% RESUMEN
% -------------------------------------------------

\textbf{Resumen}

El presente artículo realiza un análisis de los elementos necesarios para transformar una encuesta en una herramienta que nos permita recolectar información útil para una investigación. El objetivo es identificar qué características en el diseño de una encuesta influyen en la calidad y utilidad de los datos recolectados. Para ello, se hizo una revisión de documentos especializados en la metodología de la investigación y diseño de cuestionarios. Se identificó que recopilar una gran cantidad de respuestas no es sinónimo de datos útiles si las preguntas que se realizan no mantienen cohesión con el objetivo de la encuesta. Aspectos como la definición de las variables, la selección del público a entrevistar, la claridad en las preguntas y las opciones de respuesta influyen directamente en la calidad de los datos a obtener. También se encontró que una planificación deficiente puede generar información poco útil y difícil de interpretar.

\textbf{Palabras clave:} encuesta, cuestionario, información, recolección de datos, validación

% -------------------------------------------------
% DESARROLLO
% -------------------------------------------------

\textbf{Desarrollo del tema}

En teoría, realizar una encuesta es relativamente sencillo. Actualmente, cualquier persona puede crear un formulario digital, compartirlo y comenzar a acumular respuestas. El verdadero desafío aparece cuando esos datos deben utilizarse para responder a una pregunta, comprobar una teoría o tomar decisiones importantes. Tener muchas respuestas no es sinónimo de tener información útil. Si las preguntas fueron formuladas sin un propósito claro, el resultado puede ser obtener una gran cantidad de datos que aportan poco o nada al problema que se pretendía investigar.

Esta diferencia entre recolectar los datos y obtener información útil forma parte importante del diseño de encuestas. Hernández Sampieri et al. (2014) indican que la recolección de datos debe responder al planteamiento del problema. Desde esta visión, una encuesta no debería comenzar con la típica pregunta ``¿qué queremos preguntarle a las personas?'', sino ``¿qué necesitamos conocer?''. Aunque ambas preguntas parecen similares, conducen a procesos de diseño diferentes.

Definir con precisión el objetivo permite saber qué información es realmente necesaria. Por ejemplo, si una empresa desea conocer la satisfacción de sus clientes, una pregunta abierta como ``¿está satisfecho con el servicio?'' proporciona información limitada, ya que la satisfacción depende de varios factores, por ejemplo: atención, rapidez o claridad. Identificar el enfoque permite formular preguntas más específicas y posteriormente interpretar la información recolectada con mayor precisión.

La población a la que va dirigida la encuesta influye en la capacidad de obtener información útil. Una pregunta puede ser la correcta, pero puede resultar inadecuada para las personas que deben responderla. El vocabulario y el conocimiento previo que requiere la pregunta deben adaptarse a las características de los participantes.

La redacción de las preguntas representa otro punto importante. Artino et al. (2014) recomiendan utilizar preguntas claras, comprensibles y centradas en una sola idea. Las preguntas ambiguas, excesivamente complejas o que sugieren una respuesta determinada pueden introducir errores desde el momento de la recolección.

Las opciones de respuesta deben estar diseñadas con el mismo cuidado. Las preguntas cerradas facilitan la clasificación y el futuro análisis; por el lado contrario, una pregunta abierta permite identificar opiniones o situaciones que no se hayan previsto. La elección entre ambas debe basarse en el tipo de información que se desea obtener.

De esta manera, el valor de una encuesta no puede medirse únicamente por su cantidad de preguntas, participantes o gráficos generados. Su verdadero valor se encuentra en la relación entre el problema planteado, la información que se necesita conocer y la forma en que esa información será utilizada. El arte de obtener información útil consiste precisamente en formular las preguntas necesarias para transformar las respuestas de los participantes en evidencia capaz de aportar conocimiento y apoyar decisiones.

% -------------------------------------------------
% OBSERVACIONES
% -------------------------------------------------

\textbf{Observaciones y comentarios}

Una práctica útil es intentar justificar cada pregunta antes de incluirla en la encuesta. Para cada pregunta debería ser posible explicar qué se desea conocer, qué relación mantiene con el objetivo y cómo se utilizará esa respuesta en el futuro.

También se identificó que una encuesta más extensa no necesariamente proporciona mejores resultados. Reducir preguntas irrelevantes puede facilitar la participación y, al mismo tiempo, concentrar el análisis en información relacionada directamente con el problema investigado. Desde esta perspectiva, diseñar una encuesta eficiente implica también saber qué información no es necesario solicitar.

Las pruebas piloto representan otra práctica especialmente valiosa. Aunque requieren tiempo adicional, permiten detectar errores antes de que afecten a toda la recolección. Consideramos que incluso en investigaciones pequeñas debería realizarse al menos una aplicación preliminar con personas similares al público objetivo.

Como limitación, un cuestionario correctamente diseñado no garantiza por sí solo resultados representativos. Factores como la selección de la muestra, la cantidad de participantes y la tasa de respuesta también pueden influir en los resultados. Por esta razón, la calidad de las preguntas debe entenderse como una parte de un proceso de investigación más amplio.

% -------------------------------------------------
% CONCLUSIONES
% -------------------------------------------------

\textbf{Conclusiones}

\begin{enumerate}[leftmargin=*, itemsep=2pt]

    \item El propósito de una encuesta no debería limitarse a recolectar datos, sino a obtener información relacionada directamente con un problema u objetivo previamente definido.

    \item Diseñar el instrumento desde los objetivos permite seleccionar preguntas relevantes y evitar la recopilación de datos que posteriormente no tendrán una utilidad clara.

    \item La claridad y precisión de las preguntas influyen directamente en la calidad de los resultados, ya que una interpretación incorrecta durante la recolección difícilmente puede corregirse mediante el análisis posterior.

    \item La cantidad de preguntas o respuestas obtenidas no constituye por sí misma un indicador de calidad. Una encuesta breve y correctamente enfocada puede generar información más útil que un instrumento extenso sin una estructura metodológica clara.

    \item Obtener información útil mediante una encuesta requiere integrar objetivos, población, preguntas, validación y análisis como partes de un mismo proceso metodológico.

\end{enumerate}

% -------------------------------------------------
% REFERENCIAS
% -------------------------------------------------

\textbf{Referencias}

Artino, A. R., Jr., La Rochelle, J. S., Dezee, K. J., \& Gehlbach, H. (2014).
Developing questionnaires for educational research: AMEE Guide No. 87.
\textit{Medical Teacher, 36}(6), 463--474.
\url{https://doi.org/10.3109/0142159X.2014.889814}

Hernández Sampieri, R., Fernández Collado, C., \& Baptista Lucio, M. P. (2014).
\textit{Metodología de la investigación} (6.\textsuperscript{a} ed.).
McGraw-Hill Education.

% IMPORTANTE:
% La URL del repositorio debe ser la ÚLTIMA línea visible del documento.
\noindent\url{https://github.com/DaniSaurioM/ForosAcademicos-Seminario}

\textbf {Encuesta}
\noindent\url{https://forms.gle/Dmk11qtJBCMjAXTH7}

\end{document}
```
